% Source: Wald GR, physical PDF pp. 312--321 (printed pp. 299--308).
% Coverage: Section 12.1, Black Holes and the Cosmic Censor Conjecture.
% Figures/tables/footnotes: Figures 12.1--12.2; footnote 1; no equations or tables.
% Status: source-reviewed and compile-clean.
% Uncertainties: none.

\subsection{Black Holes and the Cosmic Censor Conjecture}\label{sec:12.1}
\setcounter{theorem}{0}

Our first task in the investigation of nonspherical collapse is to formulate a precise
notion of a black hole. Basically, we wish to define a black hole as a \lq\lq region of
no escape\rq\rq\ like region II of Figure~6.11, i.e., in physical terms a region of
spacetime where gravity is so strong that any particle or light ray entering that
region never can escape from it. However, this notion is \emph{not} properly captured
by defining a black hole in a spacetime \((M,g_{ab})\) to be simply a subset
\(A\subseteq M\) such that for any point \(p\in A\) we have \(J^+(p)\subseteq A\).
With that definition the causal future of any set in any spacetime would be called a
black hole. Thus, we must take much greater care to specify what portion of
spacetime the impossibility of \lq\lq escaping to\rq\rq\ should be considered grounds
for calling a region a black hole.

For asymptotically flat spacetimes, the impossibility of escaping to future null
infinity, \(\mathcal I^+\), provides an appropriate characterization of a black hole.
The crucial property of region II of Figure~6.11 which distinguishes it from, say,
the causal future of a point in Minkowski spacetime is that all of region II is confined
to \lq\lq small \(r\),\rq\rq\ i.e., region II does not extend \lq\lq out to infinity.\rq\rq\ This
is illustrated by examination of a conformal diagram of the spherical collapse
spacetime \((M,g_{ab})\) of Figure~6.11, that is, a spacetime diagram of the
unphysical spacetime \((\overline M,\widetilde g_{ab})\) associated with
\((M,g_{ab})\) (see chapter 11). This is shown in Figure~12.1. Another
representation of \(\overline M\) is given in Figure~12.2. As illustrated by these
diagrams, the causal past of future null infinity, \(J^-(\mathcal I^+)\), is
nonsingular, but it does not include the entire physical spacetime; region II is not
contained in \(J^-(\mathcal I^+)\). In contrast, for Minkowski spacetime
\(J^-(\mathcal I^+)\) includes the entire physical spacetime.

The idea that \(J^-(\mathcal I^+)\) be \lq\lq well behaved\rq\rq\ but not include the
entire spacetime leads to the following definition of a black hole. Let \((M,g_{ab})\)
be an asymptotically flat spacetime with associated unphysical spacetime
\((\overline M,\widetilde g_{ab})\). We say that \((M,g_{ab})\) is \emph{strongly
asymptotically predictable} if in the unphysical spacetime there is an open region
\(\widetilde V\subseteq\overline M\) with
\(\overline M\cap J^-(\mathcal I^+)\subseteq\widetilde V\) such that
\((\widetilde V,\widetilde g_{ab})\) is globally hyperbolic.

% Source: Wald GR, physical PDF p. 313
%         (printed p. 300).
% Coverage: Conformal diagram of the spherical-collapse spacetime.
% Figures/tables/footnotes: Figure 12.1; no tables or footnotes.
% Status: source-reviewed and compile-clean.
% Uncertainties: none.

\begingroup
\renewcommand{\thefigure}{12.1}
\begin{figure}[H]
  \centering
  \begin{tikzpicture}[x=.85cm,y=.85cm,>=Latex,line cap=round,line join=round,font=\small]
    % Cylindrical conformal compactification.
    \draw[thick] (-1.35,-2.55) -- (-1.35,2.85);
    \draw[thick] ( 1.35,-2.55) -- ( 1.35,2.85);
    \draw[thick] (-1.35,2.85) arc[start angle=180,end angle=360,x radius=1.35,y radius=.28];
    \draw[thick] (-1.35,2.85) arc[start angle=180,end angle=0,x radius=1.35,y radius=.28];
    \draw[thick] (-1.35,-2.55) arc[start angle=180,end angle=360,x radius=1.35,y radius=.28];
    \draw[dashed] (-1.35,-2.55) arc[start angle=180,end angle=0,x radius=1.35,y radius=.28];
    % The physical spacetime closes at i^0; the dashed curves are the hidden
    % halves of the same null hypersurfaces on the cylindrical diagram.
    \coordinate (bottom) at (0,-2.02);
    \coordinate (singL) at (-.28,1.78);
    \coordinate (singR) at (.28,1.78);
    \coordinate (izero) at (.88,.20);
    \draw[thick] (singR) .. controls (.92,1.42) and (1.25,.72) .. (izero);
    \draw[dashed] (singL) .. controls (-1.10,1.25) and (-1.08,.45) .. (izero);
    \draw[thick] (bottom) .. controls (.72,-1.52) and (1.23,-.62) .. (izero);
    \draw[dashed] (bottom) .. controls (-.92,-1.36) and (-.92,-.25) .. (izero);
    \node[left] at (-1.42,1.30) {\(\mathcal I^+\)};
    \node[right] at (1.35,1.18) {\(\mathcal I^+\)};
    \node[left] at (-1.40,-.98) {\(\mathcal I^-\)};
    \node[right] at (1.35,-.72) {\(\mathcal I^-\)};
    \fill (izero) circle (1.25pt) node[right] {\(i^0\)};
    % The event horizon and r=0 line meet at the regular center in the past and
    % bound the narrow black-hole region II beneath the spacelike singularity.
    \path[fill=black!10] (bottom) -- (singL) -- (singR) -- cycle;
    \draw[thick] (bottom) .. controls (-.18,-.42) and (-.26,.84) .. (singL);
    \draw[thick] (bottom) -- (singR);
    \draw[thick,decorate,decoration={zigzag,segment length=2.4pt,amplitude=.75pt}]
      (singL) -- (singR);
    \foreach \yy in {-1.55,-1.25,-.95,-.65,-.35,-.05,.25,.55,.85,1.15,1.45}
      \draw[thin] ({-.28*(\yy+2.02)/3.8},\yy) -- ({.28*(\yy+2.02)/3.8},\yy);
    \node[fill=white,inner sep=1pt] at (-.03,1.18) {II};
  \end{tikzpicture}
  \caption{A conformal diagram of the same spacetime as shown in Figures~6.11
  and~6.12. From this conformal diagram, it is apparent that region II of the physical
  spacetime lies outside of \(J^-(\mathcal I^+)\). In contrast, in Figure~11.1,
  \(J^-(\mathcal I^+)\) includes the entire physical spacetime.}
  \label{fig:12.1}
\end{figure}
\endgroup

% Source: Wald GR, physical PDF p. 313
%         (printed p. 300).
% Coverage: Another representation of the closure of the physical spacetime.
% Figures/tables/footnotes: Figure 12.2; no tables or footnotes.
% Status: source-reviewed and compile-clean.
% Uncertainties: none.

\begingroup
\renewcommand{\thefigure}{12.2}
\begin{figure}[H]
  \centering
  \begin{tikzpicture}[x=.90cm,y=.90cm,>=Latex,line cap=round,line join=round,font=\small]
    \coordinate (A) at (-1.9,-2.5);
    \coordinate (B) at (-1.9,2.5);
    \coordinate (C) at (.15,2.5);
    \coordinate (D) at (2.35,.10);
    \draw[thick] (A) -- (B);
    \draw[thick] (B) -- (C) -- (D) -- (A);
    % The event horizon is the null line from the regular center to the
    % right-hand end of the spacelike singularity.
    \draw[thick] (A) -- (C);
    \draw[thick,decorate,decoration={zigzag,segment length=2.4pt,amplitude=.75pt}]
      (B) -- (C);
    \draw[thick] (-1.90,-2.50) .. controls (-1.38,-1.45) and (-1.31,.55) .. (-1.70,1.53)
      .. controls (-1.90,2.04) and (-1.90,2.28) .. (-1.90,2.50);
    \path[fill=black!18]
      (A) -- (B) -- (-1.70,1.53) .. controls (-1.31,.55) and (-1.38,-1.45) .. (A) -- cycle;
    \node at (-.85,1.75) {II};
    \node[right] at (.67,1.91) {Event Horizon};
    \draw[->] (.57,1.84) .. controls (.27,1.82) and (.07,1.81) .. (-.15,1.77);
    \node at (1.34,1.22) {\(\mathcal I^+\)};
    \node at (1.31,-1.08) {\(\mathcal I^-\)};
    \fill (D) circle (1.25pt) node[right] {\(i^0\)};
  \end{tikzpicture}
  \caption{Another representation of the closure, \(\overline M\), of the physical
  spacetime depicted in Figure~12.1. As in Figure~12.1, the angular dimensions are
  suppressed so each point in this diagram (except those at \(r=0\) and the point
  \(i^0\)) represents a 2-sphere.}
  \label{fig:12.2}
\end{figure}
\endgroup


[Here, the closure of \(M\cap J^-(\mathcal I^+)\) is taken in the unphysical
spacetime \(\overline M\), so, in particular, \(\ii^0\in\widetilde V\). Our definition
of strong asymptotic predictability differs slightly from that given by Hawking and
Ellis 1973 since, in particular, we use a different formulation of the notion of
asymptotic flatness.] A strongly asymptotically predictable spacetime is said to
contain a \emph{black hole} if \(M\) is not contained in \(J^-(\mathcal I^+)\). The
\emph{black hole region}\footnote{The \emph{white hole region} of a strongly
asymptotically \lq\lq retrodictable\rq\rq\ spacetime is defined similarly by replacing
\(J^-(\mathcal I^+)\) with \(J^+(\mathcal I^-)\).}, \(B\), of such a spacetime is
defined to be \(B=[M-J^-(\mathcal I^+)]\), and the boundary of \(B\) in \(M\),
\(H=\partial J^-(\mathcal I^+)\cap M\), is called the \emph{event horizon}.

It should be noted that the requirement that \((\widetilde V,\widetilde g_{ab})\) be
a globally hyperbolic region of the unphysical spacetime implies that
\((M\cap\widetilde V,g_{ab})\) is a globally hyperbolic region of the physical
spacetime. Namely, according to property (1) of the definition of asymptotic
flatness given in chapter 11, we have
\[
  M=\overline M-[J^-(\ii^0)\cup J^+(\ii^0)].
\]
Hence, a Cauchy surface for \((\widetilde V,\widetilde g_{ab})\) which passes
through \(\ii^0\) will be a Cauchy surface for \((M\cap\widetilde V,g_{ab})\).
However, since \(g_{ab}\) and \(\widetilde g_{ab}=\Omega^2g_{ab}\) have the same
causal structure, this implies that \((M\cap\widetilde V,g_{ab})\) is globally
hyperbolic.

By theorem~8.3.14, we can foliate \(M\cap\widetilde V\) with Cauchy surfaces
\(\Sigma_t\). For all \(q\in M\cap\widetilde V\) and all \(\Sigma_t\) with
\(q\in J^-(\Sigma_t)\), every past directed inextendible causal curve from \(q\)
intersects \(\Sigma_t\). This can be interpreted as saying that---apart from a
possible \lq\lq initial singularity\rq\rq\ such as the white hole singularity of
Figure~6.9---no singularities are \lq\lq visible\rq\rq\ to any observer in
\([M\cap\widetilde V]\setminus[M\cap J^-(\mathcal I^+)]\). In other words, in a
strongly asymptotically predictable spacetime, all observers outside the black hole
or on the event horizon cannot \lq\lq see\rq\rq\ any singularities develop at a finite
\lq\lq time.\rq\rq\ In contrast, an asymptotically flat spacetime which fails to be
strongly asymptotically predictable is said to possess a \emph{naked singularity}.
Note that in a strongly asymptotically predictable spacetime it still is possible for
the event horizon to be singular in the sense that there may exist incomplete causal
geodesics (with respect to the physical metric, \(g_{ab}\)) in
\(M\cap J^-(\mathcal I^+)\) which are inextendible in \(M\). In particular, strong
asymptotic predictability does not exclude the possibility that the null geodesic
generators of the event horizon may be future incomplete. It is usually assumed in
discussions of black holes that the event horizon is nonsingular, but since we shall
not need any conditions beyond strong asymptotic predictability for the theorems
proven below, we shall not impose any such further conditions here.

Note also that we have defined the notion of a black hole only for strongly
asymptotically predictable spacetimes. The above definition could be given without
modification for asymptotically flat spacetimes which are not strongly asymptotically
predictable, but we choose not to do so since this case is not believed to be physically
relevant (see below) and virtually all properties of black holes derived below require
strong asymptotic predictability. The notion of a black hole also could be defined for
some non-asymptotically flat spacetimes where a suitable notion of \lq\lq infinity\rq\rq\ can
be introduced, such as spacetimes which asymptotically approach an \lq\lq open\rq\rq\
\((k=0,-1)\) Robertson--Walker universe. On the other hand, there appears to be
no natural notion of a black hole in a \lq\lq closed\rq\rq\ \((k=+1)\)
Robertson--Walker universe which recollapses to a final singularity, since there is no
natural region to which \lq\lq escape\rq\rq\ can be attempted. Of course, an
approximate notion of a black hole still exists for any region of a closed
Robertson--Walker universe that can be treated as an isolated system.


Now that we have given a precise definition of a black hole in a strongly
asymptotically predictable spacetime, we may inquire as to the physical relevance of
this definition. In the spherical case discussed in Section~6.4 above, we have seen
that gravitational collapse results in a strongly asymptotically predictable spacetime
possessing a singularity contained within a black hole. What type of spacetime
results from nonspherical collapse? First, we may invoke the singularity theorems to
conclude that---at least for sufficiently small deviations from spherical
symmetry---a spacetime singularity must occur in gravitational collapse; i.e., the
occurrence of a singularity is not merely an artifact of the assumption of exact
spherical symmetry.

To see this, we argue as follows: Starting from initial data \((\Sigma,h_{ab},K_{ab})\)
for spherical collapse, we find that trapped surfaces form in \(D^+(\Sigma)\) since
all the spheres with \(r<2M\) in the exterior Schwarzschild region are trapped
surfaces contained within \(D^+(\Sigma)\). It then follows from theorem~10.2.2 that
for all initial data sufficiently near \((\Sigma,h_{ab},K_{ab})\), trapped surfaces
also must occur in the maximal Cauchy development of these data. (Even if
departures from spherical symmetry are large, Schoen and Yau 1983 have proven that
trapped surfaces always must occur when enough matter is condensed in a small
region.) We then may appeal to theorem~9.5.3 or theorem~9.5.4 to establish the
occurrence of a spacetime singularity.

However, the existence of a spacetime singularity does not establish the existence of
a black hole, since the singularity could be naked, i.e., the spacetime could fail to be
strongly asymptotically predictable. The strongest direct evidence that naked
singularities do not occur comes from a study of linear perturbations of the
Schwarzschild spacetime. Since the metric perturbation equations (see Section~7.5)
are of the form to which theorem~10.1.2 applies, we know that an initially well
behaved perturbation on a Cauchy surface for the globally hyperbolic region
\((\widetilde V\cap M,g_{ab})\) will evolve to a nonsingular solution in
\(\widetilde V\cap M\). However, if there exists a solution of the linear perturbation
equations which \lq\lq blows up\rq\rq---i.e., is unbounded in \(\widetilde V\cap M\)---this
would indicate that the full, nonlinear equations might yield a naked singularity in
\(\widetilde V\cap M\) even for data arbitrarily close to that for Schwarzschild. On
the other hand, if all solutions of the linear perturbation equations are bounded in
\(\widetilde V\cap M\) by an appropriate norm of the initial data, this would suggest
that the solutions of the full, nonlinear equations would be nonsingular in
\(\widetilde V\cap M\) for initial data sufficiently close to that for Schwarzschild.
Studies of the behavior of solutions of the perturbation equations have shown that
the solutions are bounded in \(\widetilde V\cap M\) (see Wald 1979a). This suggests
that gravitational collapse with sufficiently small departures from spherical symmetry
will produce a black hole rather than a naked singularity. In addition, numerical
studies of linear perturbations of the Kerr black hole (see Section~12.3) have
indicated that it also is stable (Press and Teukolsky 1973).

Further support for the conclusion that gravitational collapse always produces a black
hole rather than a naked singularity comes from the fact that a number of theoretical
attempts to construct counterexamples all have failed in such a way as to suggest a
conspiracy of nature against producing naked singularities (see, e.g., Wald 1974a;
Jang and Wald 1977; Gibbons et al.\ 1983). In addition, although arguments based
on aesthetic appeal always should be viewed with suspicion, as we shall attempt to
illustrate in this chapter and in chapter 14, the study of the properties of black holes
has led to so many remarkable theoretical developments that it is very difficult to
believe that black holes are not relevant physically.

Thus, the above considerations have led to the conjecture that nature \lq\lq censors\rq\rq\
naked singularities. We may state this conjecture in physical terms as follows:

\medskip
\noindent\textsc{Cosmic Censor Conjecture} (\emph{version 1; physical formulation}).
The complete gravitational collapse of a body always results in a black hole rather
than a naked singularity; i.e., all singularities of gravitational collapse are
\lq\lq hidden\rq\rq\ within black holes, where they cannot be \lq\lq seen\rq\rq\ by distant
observers.
\medskip

The above formulation of the cosmic censor conjecture is imprecise because, among
other reasons, we have not specified what conditions the matter fields must satisfy.
Clearly, the conjecture is false without any conditions imposed on the matter fields,
since one could write down any spacetime with a naked singularity and call it a
solution of Einstein's equation by defining the stress-energy, \(T_{ab}\), to be
\((1/8\pi)G_{ab}\). Two natural conditions to impose on the matter sources are that
(1) \(T_{ab}\) satisfy an energy condition (e.g., the dominant energy condition) and
that (2) the coupled Einstein--matter field equations admit a well posed initial value
formulation. However, perfect fluids satisfy both of these conditions, but violations
of the cosmic censor conjecture with perfect fluid matter can be achieved. This can
be understood from the fact that even in Minkowski spacetime the dynamical
evolution of a perfect fluid can result in singularities such as those caused by
\lq\lq shell crossings\rq\rq\ or the formation of shocks. In the coupled
Einstein--perfect-fluid system these types of singularities still may occur, and,
indeed, because of Einstein's equation, the singularities in \(T_{ab}\) imply
singularities of the spacetime curvature as well. Thus, naked singularities can occur
in the gravitational collapse of a perfect fluid (Yodzis, Seifert, and Muller zum
Hagen 1973, 1974). However, a perfect fluid is really a macroscopic approximation
to the structure of matter rather than a fundamental description of it. Thus, the
appearance of singularities in the dynamical evolution of a fluid---particularly the
relatively \lq\lq mild\rq\rq\ types of singularities produced by shell crossings or
shocks---may be viewed as representing simply the breakdown of our macroscopic
approximation rather than the occurrence of a true, physical singularity of
gravitational collapse like the \(r=0\) singularity of the Schwarzschild solution.
This suggests that for a precise formulation of the cosmic censor conjecture we should
require the matter fields to be \lq\lq fundamental.\rq\rq\ We shall take this to mean
that the coupled Einstein--matter field equations can be put in the form of a second
order, quasilinear, diagonal, hyperbolic system (see Section~10.1), since the
fundamental classical fields which are known to accurately describe nature---namely,
gravity and electromagnetism---are of this form. Thus, we are led to the following as
one version of a mathematically precise formulation of the cosmic censor conjecture.

\medskip
\noindent\textsc{Cosmic Censor Conjecture} (\emph{version 1; precise formulation}).
Let \((\Sigma,h_{ab},K_{ab})\) be an asymptotically flat initial data set (see
chapter 11) for Einstein's equation with \((\Sigma,h_{ab})\) a complete Riemannian
manifold. Let the matter sources be such that \(T_{ab}\) satisfies the dominant
energy condition and the coupled Einstein--matter field equations are of the form
(10.1.21). In addition let the initial data for the matter fields on \(\Sigma\) satisfy
appropriate asymptotic falloff conditions at spatial infinity. Then the maximal
Cauchy evolution of these initial data (see Section~10.2) is an asymptotically flat,
strongly asymptotically predictable spacetime.
\medskip

The issue of whether the cosmic censor conjecture is correct remains the key
unresolved issue in the theory of gravitational collapse. The physical relevance of
black holes depends in large measure on the validity of this conjecture.

If the above version of the cosmic censor conjecture is correct, it is interesting to
inquire as to whether its validity, in essence, stems from a more general,
fundamental property of Einstein's equation. Penrose (1979) has suggested that this
may be the case. He has conjectured what is, in essence, a stronger version of the
cosmic censor conjecture which may be stated in physical terms as follows:

\medskip
\noindent\textsc{Cosmic Censor Conjecture} (\emph{version 2; physical formulation}).
All physically reasonable spacetimes are globally hyperbolic, i.e., apart from a
possible initial singularity (such as the \lq\lq big bang\rq\rq\ singularity) no
singularity is ever \lq\lq visible\rq\rq\ to any observer.
\medskip

This version of the cosmic censor conjecture is stronger than version 1 in that it
applies to any observer in any spacetime, not just to distant observers in
asymptotically flat spacetimes. However, version 2 does not imply version 1, since,
if, starting from asymptotically flat initial data, a singularity is formed which
propagates out to null infinity and destroys asymptotic flatness while preserving
global hyperbolicity, this would violate version 1 but not version 2.

The meaning of the second version of the cosmic censor conjecture may be
reformulated more precisely as follows. First, it certainly is not true that every
inextendible solution of Einstein's equation is globally hyperbolic. By cutting holes
in Minkowski spacetime and then making topological identifications so that not all of
the missing points can be restored, we easily may produce an inextendible,
non-globally hyperbolic spacetime. To avoid the possibility of making
\lq\lq unnecessary\rq\rq\ and \lq\lq artificial\rq\rq\ singularities of this sort, we may
give a precise formulation of the notion that all reasonable spacetimes are globally
hyperbolic as meaning that the maximal Cauchy development of nonsingular initial
data (satisfying appropriate asymptotic conditions) always yields an inextendible
spacetime. However, this property of solutions of Einstein's equation also is known
to be false. In particular, the maximal Cauchy developments of initial data for the
Taub universe (Misner 1967; Hawking and Ellis 1973) and the Kerr solution (see
Section~12.3) are known to be extendible. However, in the Taub case, the extension
violates strong causality on the Cauchy horizon, and there is good reason to believe
that if one slightly perturbs the initial data for the Taub universe in a suitable way,
one would convert the Cauchy horizon to a singularity, thereby making the maximal
Cauchy development inextendible. In the Kerr case, as discussed below, any observer
on the Cauchy horizon of the extended spacetime can \lq\lq see\rq\rq\ the entire,
noncompact initial data surface. One would expect that a suitably chosen small
perturbation of the initial data which extends out to infinity (but does not violate
asymptotic flatness) should produce an \lq\lq infinite blueshift\rq\rq\ singularity on
the Cauchy horizon. Indeed, such behavior has been explicitly demonstrated to occur
in the Reissner--Nordstrom spacetime (Chandrasekhar and Hartle 1982), whose
properties are closely analogous to those of Kerr. Thus, although some violations of
global hyperbolicity are known to occur, it appears that in these cases small
perturbations may destroy the extendibility of the maximal Cauchy development, so
that perhaps no \lq\lq generic\rq\rq\ violations are possible. Since it is difficult to
give a precise definition of the term \lq\lq generic,\rq\rq\ we formulate a precise
statement of the second version of the cosmic censor conjecture as follows (Geroch
and Horowitz 1979; Penrose 1979):

\medskip
\noindent\textsc{Cosmic Censor Conjecture} (\emph{version 2; precise formulation}).
Let \((\Sigma,h_{ab},K_{ab})\) be an initial data set for Einstein's equation, with
\((\Sigma,h_{ab})\) a complete Riemannian manifold and with the Einstein--matter
equations of the form (10.1.21) with \(T_{ab}\) satisfying the dominant energy
condition. Then, if the maximal Cauchy development of this initial data is
extendible, for each \(p\in H^+(\Sigma)\) in any extension, either strong causality
is violated at \(p\) or \(J^-(p)\cap\Sigma\) is noncompact.
\medskip

Apart from the absence of known counterexamples, there is virtually no evidence for
or against the validity of this second version of the cosmic censor conjecture.
Indeed, except for the singularity theorems (see chapter 9), very little is known
about the general, global properties of solutions of Einstein's equation.

Let us return, now, to the subject of gravitational collapse and black holes. We shall
assume the validity of the first version of the cosmic censor conjecture and shall
briefly discuss three processes by which gravitational collapse to a black hole
plausibly may occur in nature. The first process is the gravitational collapse of a
star. The evolutionary history of a star was outlined briefly in Section~6.2, where it
was concluded that if a spherical star has mass greater than about 2 solar masses, it
ultimately must undergo gravitational collapse, unless, of course, it can shed enough
mass during the course of its evolution to drop below this upper mass limit. Rotation
of the star may raise the upper mass limit, but since only one pulsar (i.e., neutron
star) has been observed to be rotating rapidly enough to possibly affect its
equilibrium structure significantly (Backer et al.\ 1982), it appears reasonable to
take the spherical upper mass limit as generally applicable. Since many stars in our
galaxy are observed to have mass greater than 2 solar masses and since the
evolutionary lifetime of massive stars is much less than the age of our Galaxy, it
would be very difficult to avoid the conclusion that many black holes have been
produced in our Galaxy by this process. However, since there are great uncertainties
as to the mass loss of stars---occurring either gradually during their evolution
(particularly in the red giant phase) or violently in a supernova explosion at the
endpoint of their evolution---there are no reliable estimates as to precisely how many
such black holes should have been produced. Perhaps the best guess can be obtained
from the observational estimate that supernovae occur in our Galaxy at the rate of
several per century. Thus at least about \(10^8\) supernovae should have occurred
during the lifetime of our Galaxy, so perhaps \(\sim10^8\) black holes may have
formed. Of course, this estimate may be far too high since many supernovae may
result in neutron stars rather than black holes, or it may be far too low since
gravitational collapse to a black hole without the violent blowing off of the outer
layers of a star may occur more frequently than supernovae. It should be noted that
black holes formed by stellar collapse must lie in the relatively narrow mass range
\(2M_\odot\lesssim M\lesssim100M_\odot\) since stars with \(M\lesssim2M_\odot\)
should not collapse, while ordinary stars with \(M\gtrsim100M_\odot\) do not exist
on account of pulsational instabilities.

A second process by which black holes may be formed in nature involves the
gravitational collapse of the entire central core of a dense cluster of stars. During the
dynamical evolution of a star cluster, gravitational encounters occasionally will
result in a large transfer of energy to a single star, which then will \lq\lq evaporate\rq\rq\
from the cluster. The remaining stars thereby lose energy and become more tightly
gravitationally bound. Thus, the central portion of a star cluster tends to become more
and more dense as the cluster evolves. Eventually a stage will be reached where the
core becomes so dense that tidal disruption of the stars will take place. It is difficult
to predict precisely what will happen after this point, but many scenarios lead to the
formation of a massive black hole (Rees 1978). (Massive black holes at the center of
star clusters also could be produced by the direct collapse of part of the gas cloud
out of which the star cluster originally formed, or by the coalescence and growth of
black holes produced by stellar collapse.) The nuclei of galaxies provide the most
likely sites for the formation of massive black holes by this process. However, again
there are no reliable estimates as to how many galactic nuclei or cores of star clusters
in other regions of a galaxy should contain a massive black hole. The mass of black
holes formed by this process could range up to a sizable fraction of the mass of the
star cluster, i.e., up to \(\sim10^{10}M_\odot\) for a black hole in the nucleus of a
large galaxy.

A third, much more speculative, process by which black holes may have been
produced is by the gravitational collapse of regions of enhanced density in the early
universe. As discussed in chapter 5, the universe appears to be homogeneous and
isotropic to an excellent approximation on large distance scales. However, the matter
distribution we observe in the present universe certainly is not exactly homogeneous,
so some initial inhomogeneities must have been present in the early universe. Some
information about the initial spectrum of density inhomogeneities on galactic mass
scales can be obtained by statistical studies of the present clustering behavior of
galaxies (Peebles 1980), but very little is known about initial inhomogeneities on
smaller scales. However, if sufficiently large inhomogeneities in the density were
present in the early universe, the regions of enhanced density could collapse directly
to a black hole rather than expand with the rest of the universe. Thus, large numbers
of \lq\lq primordial black holes\rq\rq\ could have been produced in this way. Again,
however, there are no reliable predictions of how many---if any---primordial black
holes exist in our universe. Indeed, the best limits on the density inhomogeneities on
small mass scales come from the requirement that they not overproduce small black
holes.

An important feature of primordial black holes is that they could have been produced
at any mass scale, including masses much smaller than a solar mass. No process in
the present universe could produce such small black holes. For example, a
Schwarzschild black hole of mass equal to that of the Earth,
\(M_\oplus=6\times10^{27}\,\mathrm{g}\), has \(r_s=2GM_\oplus/c^2\sim1\,\mathrm{cm}\).
Thus, to convert the Earth into a black hole, we would, presumably, have to compress
it by nongravitational forces down to a radius of order \(r_s\) (and thus a density
\(\sim10^{27}\,\mathrm{g/cm^3}\)) before self-gravitation would result in its collapse
to a black hole. On the other hand, at \(T\sim10^{-11}\,\mathrm{s}\) after the big bang,
the density, \(\rho\), of the universe was \(\sim10^{27}\,\mathrm{g/cm^3}\). A density
fluctuation with \(\delta\rho\sim\rho\) over a scale of \(1\,\mathrm{cm}\) at this time
could have resulted in the formation of a primordial black hole with mass
\(M_\oplus\). Thus, very soon after the big bang, black holes of small mass may have
been produced.

How might black holes produced by the above three processes be detected? An
important point to recognize is that black holes are extremely small objects. A black
hole of one solar mass has a Schwarzschild radius of only 3 km; even a black hole
of \(10^{10}M_\odot\) in the nucleus of a large galaxy would have
\(r_s\sim3\times10^{10}\,\mathrm{km}\sim3\times10^{-3}\) light-years. When added to
the fact that black holes are \lq\lq black\rq\rq\ (see problem 1), it is clear that direct
detection of a black hole would be extremely difficult. The most promising
possibility for indirect detection of a black hole comes from the fact that matter
which accretes onto it will heat up and emit electromagnetic radiation before entering
the black hole. For black holes formed by stellar collapse, the best opportunity for
such accretion occurs if the black hole is in a close binary orbit with a star, so that
matter can flow from the star toward the black hole. In this case, the matter would be
expected to slowly spiral into the black hole, thereby forming an accretion disk around
the black hole. Viscous heating in the accretion disk could result in the production of
X-rays. A number of X-ray sources with an ordinary (i.e., uncollapsed) star in close
binary orbit around an unseen (at optical frequencies) companion have been
discovered, but accretion onto a neutron star (and possibly even a white dwarf) also
could produce such X-rays, so one cannot conclude that these systems must contain a
black hole. However, in the case of Cygnus X-1, the properties of the binary orbit
yield a rather firm lower mass limit for the unseen companion of \(\sim9M_\odot\)
(Paczynski 1974). This is far above the upper mass limit for neutron stars and white
dwarfs, so there is very strong reason to believe that the unseen companion of the
Cygnus X-1 system is a black hole. Very recently, the X-ray source LMC X-3 also
has been determined to consist of a binary system, with the mass of the unseen
companion well above the white dwarf and neutron star limits (Cowley et al.\ 1983;
Paczynski 1983).

A massive black hole at the center of a star cluster or galactic nucleus could produce
an observable effect by altering the equilibrium distribution of stars in the central
portion of the cluster, causing more stars to be \lq\lq drawn in\rq\rq\ toward the center.
Thus, if a massive black hole were present in a star cluster, one would expect to see a
small brightness enhancement very near the center of the cluster. In addition, one
would expect to see an increase in the average velocity of stars very near the center.
Exactly such a brightness enhancement and increased \lq\lq velocity dispersion\rq\rq\
have been observed at the center of the galaxy M87 (Young et al.\ 1978; Sargent et
al.\ 1978), thus providing strong evidence for the presence there of a black hole of
mass \(\sim5\times10^9M_\odot\). Furthermore, the galaxy M87 is well known for the
existence of a jet of highly energetic particles emanating from the center of the
galaxy. Thus, there is a strong suggestion that a massive black hole may be involved
in the mechanism which produces this jet, although at present there is no fully
satisfactory theory of exactly how a black hole (or any other object) could produce
such a jet. Since similar jets occur in other active galaxies and in quasars, massive
black holes may be present in these bodies as well. Indeed, as already mentioned at
the beginning of this book, the discovery of quasars in the early 1960s and the
inability to explain their energy source without the presence of strong gravitational
fields provided a great stimulus to the development of the theory of gravitational
collapse and black holes.

Finally, we mention a possible means of detecting primordial black holes of very low
mass. Classically such black holes would produce virtually no observable effects
unless they were sufficiently numerous to provide a cosmologically significant
contribution to the mass density of the universe or unless one of them happened to
strike the Earth. However, as we shall see in chapter 14, particle creation occurs near
black holes and is appreciable for black holes of very small mass. The effects of
numerous primordial black holes with \(M\lesssim10^{15}\,\mathrm{g}\) could
significantly contribute to the \(\gamma\)-ray background (Page and Hawking 1976).
However, no such contribution is seen, so one only may place an upper limit on the
number of such small black holes by this means.

